% !TEX root = ../dissertationOG.tex
% Chapter 6 revised to reflect the completed studies in Chapters 3--5.
% Existing chapter and section labels are preserved for cross-references.

\clearpage
\chapter{Limitations and Future Directions}\label{CH6_LIMITATION}

The preceding chapters demonstrate three complementary uses of physics-based particle simulation: controlling deformable motion through optimized internal stresses (Chapter~\ref{CH3_SHAPEMORPHING}), incorporating image-space feedback into that control (Chapter~\ref{CH4_PhysMorphGS}), and generating and evaluating physically grounded robotic cutting demonstrations (Chapter~\ref{CH5_CUTTINGSIM}). Although these studies share the Material Point Method (MPM) as a modeling foundation, they address different objectives and do not constitute a single end-to-end differentiable system. This chapter summarizes the limitations of each study and identifies directions in which their methods and evaluations can be strengthened.

\section{Limitations of Physics-Based Shape Morphing}
\label{subsec:lim-mpm}

\subsection{Time Integration and Optimization Cost}
The differentiable MLS-MPM method in Chapter~\ref{CH3_SHAPEMORPHING} uses explicit time stepping and relatively small simulation time steps; the reported morphing examples typically use $\Delta t=1/120$\,s. This numerical time step should not be confused with a measured real-time execution rate. Gradient-based control requires repeated forward simulations and backward passes, so the wall-clock cost grows with the number of simulation steps, control intervals, and optimization iterations. Damping and numerical stabilization improve robustness but can also change the temporal character of the resulting motion. More efficient integration and control optimization would be needed for applications requiring rapid interactive updates.

\subsection{Limited Optimization Horizons}
Differentiating through many MPM time steps requires retaining or reconstructing intermediate simulation states. Long unrolled trajectories can therefore be costly in memory and susceptible to poorly conditioned or attenuated gradients. Chapter~\ref{CH3_SHAPEMORPHING} addresses this problem by optimizing shorter simulation networks and chaining their final states. This permits longer morphing sequences, but it does not optimize all control variables over an arbitrarily long sequence simultaneously. Consequently, local improvements need not yield a globally optimal long-horizon trajectory.

\subsection{Dependence on the Objective and Smoothing Parameters}
The correspondence-free, log-based nodal-mass objective reduces the mass ejection observed with a direct nodal-mass discrepancy. Nevertheless, the need for a mass penalty and carefully selected optimizer settings shows that optimization behavior is sensitive to the objective's scale and shape. Matching grid-mass distributions also does not guarantee an equally accurate rendered boundary or recovery of every thin feature. In the reported examples, deformation-gradient smoothing factors are selected within approximately $\gamma=0.95$--$0.98$ to suppress oscillatory artifacts. Such smoothing is useful for visual stability but can reduce sensitivity to rapid local deformation, and its preferred strength may change with geometry, material stiffness, and discretization.

\section{Limitations of Render-Guided Volumetric Morphing}
\label{subsec:lim-bridge}

\subsection{The Render Term's Authority after Arrival}
In Chapter~\ref{CH4_PhysMorphGS} the render weight is set in every window so that the render term holds a fixed share of the first gradient. While the body travels to the target this gives the render term real influence: it redirects material well beyond the run-to-run noise. Once the body has arrived, however, the physics gradient becomes small, the render weight follows it, and the render term only polishes the outline. The end is consistently but only slightly closer to the target than without the render term. The physics objective reads only the positions of the particles, so it cannot tell which material ends where, nor relief finer than the three-dimensional target resolves, and these are exactly the things a picture could decide. Holding the render weight fixed lets the render term act on them, but the body then keeps moving at the end of the run. A weight that keeps the render term's authority where it has something to decide and still lets the body come to rest has yet to be found.

\subsection{Surface Detail below the Grid}
Much of what a picture asks for lies below one MPM cell. The rigid-free offset of the outer particle layer reaches this scale, but it is also the main source of the few particles that briefly separate from the body while it arrives. The interpolation kernel averages the grid velocity over about two cells, so gaps of the target that are narrower than this open late and can tear into beads when they do. A finer cell opens some of these gaps but costs time and leaves holes in the early surface, so the cell size remains a trade-off to be settled by further comparison.

\subsection{Rendering Representation versus Physical State}
The render terms read a display surface built from the particles, the exterior of their own field covered by surface discs; this surface carries no mass of its own and is not a second physical discretization. An improved silhouette therefore cannot certify the accuracy of internal stress, contact response or material parameters, and part of the render term's lead on its own measure comes from a smoother arrangement of the outer layer rather than a different shape, which is why Chapter~\ref{CH4_PhysMorphGS} reads shape against the exact target mesh. The weak ambient drag does not conserve angular momentum, and most of the delivered frames come from a long, slow tail after the body has arrived. Joint evaluation of rendering quality and physical behavior remains important.

\section{Limitations of Physics-Grounded Robotic Cutting}
\label{subsec:lim-cutting}

\subsection{Geometric Generalization and Sequential Actions}
The CulinaryCut experiments in Chapter~\ref{CH5_CUTTINGSIM} expose a \emph{geometry gap}: successful cutting on familiar objects or instructions does not ensure accurate placement under changes in object scale, target identity, directional reference, or requested cut ratio. Multi-object scenes and held-out food categories reduce performance, and sequential split cuts remain particularly difficult because the next cutting location must be recomputed after the geometry changes. These observations indicate that short-horizon action prediction does not by itself provide reliable object-centric or long-horizon reasoning.

\subsection{Approximate Cutting Mechanics and Material Calibration}
IntSim combines MLS-MPM, signed-distance representations of the knife and board, contact and friction handling, and a localized damage mechanism. Its damage thresholds, blade-proximity bands, and resolution-dependent scalings are practical modeling choices rather than a general guarantee of resolution-independent fracture. Real food materials may exhibit heterogeneity, rate dependence, adhesion, and fracture behavior beyond the chosen constitutive and damage approximations. Material parameters are initialized from available information and adjusted where necessary for numerical robustness; the work does not establish independent constitutive calibration for all 15 food categories. More extensive measurements would be necessary before claiming material-independent predictive force fidelity.

\subsection{Limits of Physical Validation and Sim-to-Real Transfer}
The reported force validation compares real and simulated banana cutting, including an aggregate comparison of 20 real and 30 simulated trials. These measurements provide useful evidence for a representative contact interaction, but cannot establish force accuracy across every food, knife motion, or material condition. The real-robot orange-cutting comparison is also preliminary: the two RDT training conditions yield 1/10 and 3/10 successful trials, respectively. This sample is insufficient to infer robust transfer across objects, embodiments, or deployment settings. Furthermore, CulinaryCut's approximately 1.575 million augmented observation--instruction instances are derived from 21,000 underlying physical trajectories; visual and language augmentation should not be counted as independent physics rollouts.

The evaluated VLA architectures do not receive simulated or measured contact forces as policy inputs. Consequently, the observed improvement from physics-grounded training data supports \emph{data-level physical grounding}, not the acquisition of explicit online force feedback. Likewise, the execution-time velocity guardrail is an external safety mechanism rather than a learned force-conditioned controller. Incorporating force sensing and feedback into learned policies remains a separate research direction.

\section{Shared Challenges in Efficiency and Reproducibility}
\label{subsec:lim-scope}

The three studies place different demands on computation. The shape-morphing experiments in Chapter~\ref{CH3_SHAPEMORPHING} report a CPU-parallel MPM implementation, Chapter~\ref{CH4_PhysMorphGS} runs on one GPU and adds Sinkhorn transport solves, surface searches, rendering, and an adjoint through every optimization window, and the IntSim implementation in Chapter~\ref{CH5_CUTTINGSIM} uses GPU-based Taichi kernels for material and contact computations. A limitation of one implementation should therefore not be generalized to all three systems. Reproducible comparisons require separate reporting of simulation resolution, physical simulation duration, wall-clock runtime, optimization budget, and hardware.

Results can also depend on material parameters, contact thresholds, particle and grid resolution, sampling seeds, and optimizer schedules. Future evaluations should disclose these settings, distinguish numerical robustness from physically measured agreement, and report variation across multiple trials. In particular, visual quality, geometric task success, material separation, and measured interaction forces should remain distinct evaluation criteria rather than being combined into an unsupported claim of universal physical fidelity.

\section{Future Research Directions}
\label{subsec:future}

\subsection{More Efficient and Stable Simulation-Based Optimization}
A promising direction is to investigate implicit or mixed time integration, adaptive time stepping, and more effective control parameterizations for large-deformation morphing. Checkpointing, recomputation, and memory-aware adjoint methods could reduce the cost of long differentiable simulations. These improvements should be evaluated against both convergence quality and actual computational expense, rather than assuming that larger numerical time steps necessarily make inverse optimization faster.

\subsection{Lasting Rendering Feedback and Finer Surface Control}
For render-guided morphing, the next step is a render weight that keeps its authority where the picture has something to decide, such as a color layout, lettering or relief finer than the three-dimensional target, and still lets the body come to rest. Further directions are a symmetric link between neighboring steps of the layer offset that removes transient floating particles without moving the center of mass, a choice of the MPM cell that opens narrow gaps at an acceptable cost, a drag that damps the whole APIC momentum, and targets given only by pictures, where the concavities that only the images contain must be carved from a visual hull.

\subsection{Measured Material Properties and Broader Robotic Validation}
Future work should identify constitutive, contact, and damage parameters from controlled measurements across additional foods, tools, and cutting styles. Repeated real-robot trials should quantify force--time uncertainty and evaluate cluttered scenes, sequential cuts, and active object stabilization. Force-conditioned policies should be investigated as a separate extension.

\subsection{Bridging Simulation Objectives and Downstream Task Success}
Across the three studies, an open question is how to connect numerical simulation quality to what ultimately matters for a particular application. In animation, this includes deformation control and rendered temporal coherence; in robotic manipulation, it includes correct geometry grounding, physical separation, and safe interaction. Benchmarks that report these outcomes alongside discretization error, material sensitivity, and computational cost would help determine when increased simulation fidelity produces meaningful improvements downstream.

\section{Concluding Remarks}
The dissertation demonstrates that MPM-based simulation can support physically directed shape transformations, differentiable rendering-guided control, and physics-grounded data generation for robotic cutting. Each application benefits from a common representation of deformable material motion, while requiring different optimization, rendering, or evaluation components. The remaining challenges concern both numerical fidelity and the connection between simulation outputs and real application outcomes. Addressing them will require more efficient optimization, more reliable surface and contact models, and broader validation against observed physical behavior.
